\section{Introduction}
\label{sec:intro}

\subsection{Problem statement}
Before granting a loan, a lender must estimate how likely the applicant is to default, that is, to stop repaying.
Approving a future defaulter loses the principal, while rejecting a reliable applicant only loses the interest margin, so the two kinds of error have very different costs.
This report studies the problem as binary classification on the public Credit Risk Dataset \cite{kaggle_credit}.
Each of its \num{32581} records describes one applicant (age, income, employment length, home ownership, credit history) and one loan (purpose, grade, amount, interest rate, loan-to-income ratio).
The target \ident{loan_status} equals 1 when the loan defaulted.

The dataset meets the constraints of the assignment for tabular data.
It contains missing values (in \ident{person_emp_length} and \ident{loan_int_rate}), which calls for imputation.
It contains four categorical variables, which call for encoding.
It is large enough for reliable cross-validation and a held-out test set of \num{6484} loans.

\subsection{Objectives}
Following the assignment specification, the project pursues four objectives.
\begin{enumerate}
  \item \textbf{Traditional pipeline (mandatory).} Build a complete pipeline with exploratory data analysis, preprocessing, feature extraction and selection, model training and evaluation, where every step is configurable. Then measure the effect of each choice with controlled experiments.
  \item \textbf{Model comparison.} Tune and compare eight classical classifiers under one evaluation protocol designed for an imbalanced, cost-asymmetric problem.
  \item \textbf{Deep learning pipeline (bonus).} Train neural networks that learn their own representation of the tabular data (an entity-embedding MLP and TabNet), save the learned features, and compare them with the traditional pipeline.
  \item \textbf{Bayesian network (talented-engineering extension).} Implement a Bayesian network from scratch, with structure and parameter learning, exact inference (variable elimination) and approximate inference (rejection sampling, likelihood weighting, Gibbs sampling). Then apply it to the same prediction problem and to probabilistic reasoning about the causes of default.
\end{enumerate}

\subsection{Contributions}
\begin{itemize}
  \item A reusable Python package (\ident{modules/creditrisk}) that builds any pipeline variant from a declarative configuration. It is driven by a Google Colab notebook that runs end to end with \emph{Run all}.
  \item Six ablation studies on imputation, scaling, outliers and feature engineering, encoding, dimensionality reduction, and class imbalance, followed by grid-searched models with F2-tuned decision thresholds.
  \item A deep learning comparison with seven MLP variants and TabNet, and an analysis of the learned embeddings.
  \item A from-scratch Bayesian-network library, verified by unit tests against brute-force enumeration and scikit-learn's \ident{CategoricalNB}. It is used to compare four network structures and three sampling algorithms.
  \item Saved feature sets (\ident{.npy} and \ident{.h5}) for every pipeline, and a public GitHub repository.
\end{itemize}

\subsection{Deliverables and reproducibility}
The submission follows the required layout: \ident{notebooks/}, \ident{modules/}, \ident{reports/} and \ident{features/}.
\begin{itemize}[nosep]
  \item Source code: \url{https://github.com/datduongabc/credit-risk-classification}
  \item Colab notebook: \url{https://colab.research.google.com/github/datduongabc/credit-risk-classification/blob/main/notebooks/credit_risk_pipeline.ipynb}
\end{itemize}
The notebook clones the repository, installs the missing libraries and downloads the dataset from a public link, without mounting any personal storage.
Every number in this report is read from files that the notebook generates (\ident{reports/generated/}).
All randomness is seeded, so a re-run reproduces the tables; the deep models, whose results vary slightly across hardware and PyTorch builds, are reported as the mean and standard deviation over five seeds, so a run on another machine lands within the reported spread.

\subsection{Organization of the report}
\Cref{sec:eda} analyses the data and derives the preprocessing decisions.
\Cref{sec:method} presents the evaluation protocol and the three modelling approaches.
\Cref{sec:experiments} reports the experiments.
\Cref{sec:discussion} discusses strengths, limitations and threats to validity.
\Cref{sec:conclusion} concludes.
The appendix documents the code structure, the tests and the saved features.
